% GENERATED FILE. Edit canonical lessons, then run npm run generate:books.
% canonical-source-hash: fnv1a64:f13da615f495e06c
% canonical-lessons: ZH-C405-zhida, ZH-C405-jinru, ZH-W405-ru-enter, ZH-C405-churu, ZH-C405-mobanche, ZH-C405-tingkao

\chapter{Through the Gates}
\label{ch:zh-through-the-gates}
\hlchaptermodality{405}

\begin{chapteropening}
\textbf{By the end of this chapter:} \emph{I can say a train goes straight through, go in and out of the station, and talk about the last train.}

Complete the last lesson of chapter 405: I can say a train goes straight through, go in and out of the station, and talk about the last train.
\end{chapteropening}

\section[zhídá]{zhídá — to go direct; to go without changing}

\label{lesson:ZH-C405-zhida}



\begin{quote}
Before the new one: say the Mandarin for a sleeping berth, then the Mandarin for a hard seat.
\end{quote}

\subsection*{You'll want to know: zhídá}
\emph{zhídá} — \textquotedblleft{}to go direct; to go without changing\textquotedblright{}.

\begin{quote}
\emph{zuò gāotiě kěyǐ zhídá Běijīng} — The high-speed train takes you straight to Beijing.
\end{quote}

Its characters come later; for now it is a word for the ear.

\begin{cousinweb}
\emph{zhídá} is built from \emph{zhí} 'straight' and \emph{dá} 'to get through', both inherited from Middle Chinese, \emph{drik} and \emph{dat} (Guangyun, in Baxter's transcription), which glosses the first 'upright' and the second 'to reach through'.
\end{cousinweb}

\begin{grammarlens}[title={What you've built}]
One more word for getting around.
\end{grammarlens}

\subsection*{Your turn}
\begin{itemize}
\raggedright
  \item \emph{Say it:} \emph{zhídá}
  \item \emph{Say it:} \emph{zhídá}, once more
  \item \emph{Say it:} \emph{zuò gāotiě kěyǐ zhídá Běijīng} — The high-speed train takes you straight to Beijing.
  \item \emph{Recall:} say the Mandarin for a sleeping berth, then the Mandarin for a hard seat, then say \emph{zhídá} again
\end{itemize}

\begin{tcolorbox}[breakable,colback=teal!4,colframe=teal!35!black,title={Before you move on}]
What does zhídá mean? (\textquotedblleft{}To go direct; to go without changing\textquotedblright{}.) Say it once more.
\end{tcolorbox}
\section[\texorpdfstring{\zh{进入} (jìnrù)}{进入 (jìnrù)}]{\zh{进入} (jìnrù) — to go into}

\label{lesson:ZH-C405-jinru}



\begin{quote}
Before the new one: say the Mandarin for a hard seat, then the Mandarin for to go without changing.
\end{quote}

\subsection*{You'll want to know: \zh{进入}}
\textbf{\zh{进入}} — \emph{jìnrù} — \textquotedblleft{}to go into\textquotedblright{}.

\begin{quote}
\emph{wǎng qián zǒu jiù kěyǐ jìnrù chēzhàn} — Go straight on and you can get into the station.
\end{quote}

The new character, in pieces: \zh{入} — two strokes: a short stroke falling to the left, then a long stroke falling to the right that starts higher, above the top of the first. The next lesson puts it in your hand.

\begin{cousinweb}
\textbf{\zh{进入}} (\emph{jìnrù}) is built from \zh{进} 'to go forward' and \zh{入} 'to go in', both inherited from Middle Chinese, \emph{tsinH} and \emph{nyip} (Guangyun, in Baxter's transcription), which glosses the first 'forward' and the second 'within' and 'to take in'.
\end{cousinweb}

\begin{grammarlens}[title={What you've built}]
One more word for getting around.
\end{grammarlens}

\subsection*{Your turn}
\begin{itemize}
\raggedright
  \item \emph{Say it:} \emph{jìnrù}
  \item \emph{Say it:} \emph{jìnrù}, once more
  \item \emph{Say it:} \emph{wǎng qián zǒu jiù kěyǐ jìnrù chēzhàn} — Go straight on and you can get into the station.
  \item \emph{Recall:} say the Mandarin for a hard seat, then the Mandarin for to go without changing, then say \emph{jìnrù} again
\end{itemize}

\begin{tcolorbox}[breakable,colback=teal!4,colframe=teal!35!black,title={Before you move on}]
What does \zh{进入} mean? (\textquotedblleft{}To go into\textquotedblright{}.) Say it once more.
\end{tcolorbox}
\section[\texorpdfstring{\zh{入} (rù)}{入 (rù)}]{\zh{入} — two strokes, the right one on top}

\label{lesson:ZH-W405-ru-enter}



\begin{quote}
The word first, then the character inside it.

\begin{itemize}
\raggedright
  \item \emph{Recall:} say \emph{jìnrù} — \textbf{R1}, one lesson back, before it decays
\end{itemize}

\begin{quote}
\zh{人} \zh{进}
\end{quote}

\emph{Read it:} each one aloud, with nothing said first — \emph{rén}, \emph{jìn}
\end{quote}

\subsection*{Script — two strokes, the right one on top}
\hlblockfigure{\detokenize{figures/ZH-W405-ru-enter-filmstrip.pdf}}{How \zh{入} is written, stroke by stroke}

\begin{quote}
\zh{入}
\end{quote}

\textbf{One piece, two strokes.}

\textbf{First, a short stroke falling to the left}, from the middle down to the lower left.

\textbf{Then a long stroke falling to the right.} It starts higher, above the top of the first, and sweeps down to the lower right.

In \zh{人} the left stroke is the long one and the right stroke leans on it; in \zh{入} it is the other way round.

\subsection*{Your turn}
\begin{itemize}
\raggedright
  \item \emph{Write it:} \zh{入}, all two, with the model in front of you
  \item \emph{Cover:} write it again from memory, and count the pieces, not the strokes
\end{itemize}

\begin{tcolorbox}[breakable,colback=teal!4,colframe=teal!35!black,title={Before you move on}]
How many strokes in \zh{入}? (\textbf{Two}.) Which one starts higher? (\textbf{The second, falling to the right}.) How is it unlike \zh{人}? (\textbf{In \zh{入} the right stroke is the long one}.)

Source: \href{https://raw.githubusercontent.com/chanind/hanzi-writer-data/68d10a4b21150cae5e1ebbd223eed289cf32d90c/data/\%E5\%85\%A5.json}{Hanzi Writer Data \zh{入}, pinned stroke paths and medians}.
\end{tcolorbox}
\section[\texorpdfstring{\zh{出入} (chūrù)}{出入 (chūrù)}]{\zh{出入} (chūrù) — to go in and out}

\label{lesson:ZH-C405-churu}



\begin{quote}
Before the new one: say the Mandarin for to go without changing, then the Mandarin for to go into a place.
\end{quote}

\subsection*{You'll want to know: \zh{出入}}
\textbf{\zh{出入}} — \emph{chūrù} — \textquotedblleft{}to go in and out\textquotedblright{}.

\begin{quote}
\textbf{\zh{旅客从这个门出入}} \emph{lǚkè cóng zhège mén chūrù} — Passengers go in and out through this gate.
\end{quote}

\begin{cousinweb}
\textbf{\zh{出入}} (\emph{chūrù}) is built from \zh{出} 'to go out' and \zh{入} 'to go in', both inherited from Middle Chinese, \emph{tsyhwit} and \emph{nyip} (Guangyun, in Baxter's transcription).
\end{cousinweb}

\begin{grammarlens}[title={What you've built}]
One more word for getting around.
\end{grammarlens}

\subsection*{Your turn}
\begin{itemize}
\raggedright
  \item \emph{Say it:} \emph{chūrù}
  \item \emph{Say it:} \emph{chūrù}, once more
  \item \emph{Say it:} \emph{lǚkè cóng zhège mén chūrù} — Passengers go in and out through this gate.
  \item \emph{Recall:} say the Mandarin for to go without changing, then the Mandarin for to go into a place, then say \emph{chūrù} again
\end{itemize}

\begin{tcolorbox}[breakable,colback=teal!4,colframe=teal!35!black,title={Before you move on}]
What does \zh{出入} mean? (\textquotedblleft{}To go in and out\textquotedblright{}.) Say it once more.
\end{tcolorbox}
\section[mòbānchē]{mòbānchē — last bus or train}

\label{lesson:ZH-C405-mobanche}



\begin{quote}
Before the new one: say the Mandarin for to go into a place, then the Mandarin for to go in and out.
\end{quote}

\subsection*{You'll want to know: mòbānchē}
\emph{mòbānchē} — \textquotedblleft{}last bus or train\textquotedblright{}.

\begin{quote}
\emph{wǒmen zuò mòbānchē qù jīchǎng} — We're taking the last train to the airport.
\end{quote}

Its characters come later; for now it is a word for the ear.

\begin{cousinweb}
\emph{mòbānchē} is built from \emph{mò} 'last', \emph{bān} 'a shift' and \zh{车} 'vehicle', all inherited from Middle Chinese, \emph{mat}, \emph{paen} and \emph{tsyhae} (Guangyun, in Baxter's transcription). The dictionary glosses the first as the tip of a tree and 'the end', and the second as dividing jade tokens, so the sense 'a shift, a run' is not given there.
\end{cousinweb}

\begin{grammarlens}[title={What you've built}]
One more word for getting around.
\end{grammarlens}

\subsection*{Your turn}
\begin{itemize}
\raggedright
  \item \emph{Say it:} \emph{mòbānchē}
  \item \emph{Say it:} \emph{mòbānchē}, once more
  \item \emph{Say it:} \emph{wǒmen zuò mòbānchē qù jīchǎng} — We're taking the last train to the airport.
  \item \emph{Recall:} say the Mandarin for to go into a place, then the Mandarin for to go in and out, then say \emph{mòbānchē} again
\end{itemize}

\begin{tcolorbox}[breakable,colback=teal!4,colframe=teal!35!black,title={Before you move on}]
What does mòbānchē mean? (\textquotedblleft{}Last bus or train\textquotedblright{}.) Say it once more.
\end{tcolorbox}
\section[tíngkào]{tíngkào — to call at a station or port}

\label{lesson:ZH-C405-tingkao}



\begin{quote}
Before the new one: say the Mandarin for to go in and out, then the Mandarin for the last bus or train.
\end{quote}

\subsection*{You'll want to know: tíngkào}
\emph{tíngkào} — \textquotedblleft{}to call at a station or port\textquotedblright{}.

\begin{quote}
\emph{gāotiě zài xiǎo zhàn bù tíngkào} — The high-speed train doesn't stop at small stations.
\end{quote}

Its characters come later; for now it is a word for the ear.

\begin{cousinweb}
\emph{tíngkào} is built from \emph{tíng} 'to stop' and \emph{kào} 'to lean against', both inherited from Middle Chinese, \emph{deng} and \emph{khawH} (Guangyun, in Baxter's transcription), which glosses the first 'to rest, to stop' and the second 'to go against each other'; the sense 'to lean on' is not given there.
\end{cousinweb}

\begin{grammarlens}[title={What you've built}]
One more word for getting around.
\end{grammarlens}

\subsection*{Your turn}
\begin{itemize}
\raggedright
  \item \emph{Say it:} \emph{tíngkào}
  \item \emph{Say it:} \emph{tíngkào}, once more
  \item \emph{Say it:} \emph{gāotiě zài xiǎo zhàn bù tíngkào} — The high-speed train doesn't stop at small stations.
  \item \emph{Recall:} say the Mandarin for to go in and out, then the Mandarin for the last bus or train, then say \emph{tíngkào} again
\end{itemize}

\begin{tcolorbox}[breakable,colback=teal!4,colframe=teal!35!black,title={Before you move on}]
What does tíngkào mean? (\textquotedblleft{}To call at a station or port\textquotedblright{}.) Say it once more.
\end{tcolorbox}
